Quá trình khảo sát đã thực hiện thành công việc mô phỏng và phân tích các hiện tượng dao động phi tuyến thông qua thuật toán Runge-Kutta bậc 4 (RK4). Sự tương đồng giữa kết quả tính toán số và các giá trị lý thuyết đối chứng đã khẳng định tính chính xác, độ tin cậy của mô hình toán học và phương pháp giải được trình bày trong bài báo này.

Dựa trên các kết quả mô phỏng, chúng tôi rút ra các kết luận chính sau:
\begin{itemize}
    \item Về ma sát: Thuật toán đã tái hiện chính xác ba trạng thái cản (dưới cản, tới hạn và quá cản). Đặc biệt, hệ số cản tới hạn được xác định là ngưỡng tối ưu để triệt tiêu dao động nhanh nhất cho hệ thống.
    \item Về cộng hưởng: Hệ thống bộc lộ phản ứng biên độ cực đại khi tần số ngoại lực tiệm cận tần số riêng. Việc mô phỏng chi tiết vùng cộng hưởng giúp đưa ra các cảnh báo quan trọng về giới hạn bền của cấu trúc cơ khí.
    \item Về tính phi tuyến: Sự thay đổi bậc thế năng ($p=4$) và tham số $\alpha$ đã làm thay đổi bản chất của hệ: chu kỳ dao động không còn là hằng số mà phụ thuộc vào biên độ, đồng thời quỹ đạo dao động trở nên bất đối xứng.
\end{itemize}

Tóm lại, mặc dù có những sai lệch nhỏ so với thực tế do các yếu tố môi trường, các số liệu tính toán này vẫn là cơ sở khoa học tin cậy cho việc dự báo, thiết kế và tối ưu hóa các thông số vận hành nhằm đảm bảo tính ổn định cao nhất cho các hệ thống dao động thực tế.
